\section{The metric adapted to thermal relaxation}\label{sec:metric}

Fix a minimizing pair from Proposition~\ref{prop:regularized-minimum}.
All comparisons in this section concern these states.  We use the centered
observable spaces, modular frequencies $v_s$, and maps $L_i$ introduced in
Section~\ref{sec:minimization}; sums over $i$ run over $1,2$.
For a positive scalar weight $X_s(v_s)$, write
\[
 \|Z\|_{X_s,s}^2=\langle Z,X_sZ\rangle_s.
\]
We say that a family of weights satisfies comparison when
\begin{equation}\label{eq:weight-comparison}
 \sum_i w_i\|L_iZ\|_{X_i,i}^2\le \|Z\|_{X_0,0}^2.
\end{equation}
Initially, $Z$ is a finite eigenblock matrix centered by a scalar multiple
of the identity.  Such tests have finitely many frequencies, including zero;
they need not have finite matrix support after centering.

\subsection{Four elementary comparisons}
Define, with continuous values at zero,
\[
 b_s(v)=\frac v2\coth\frac{h_sv}{2},\qquad
 R_{s,\pm}(v)=b_s(v)\pm\frac v2,
 \qquad b_s(0)=\frac1{h_s}.
\]
If $\rho_s=\operatorname{diag}(p_l)$, then
\begin{equation}\label{eq:arithmetic-weights}
 h_s\ell(p_l,p_r)R_{s,+}(v_s(l,r))=p_l,
 \qquad
 h_s\ell(p_l,p_r)R_{s,-}(v_s(l,r))=p_r.
\end{equation}
Thus the corresponding squared norms of a centered observable are
$\Tr\rho_sZZ^\dagger$ and $\Tr\rho_sZ^\dagger Z$.

Both families $R_+$ and $R_-$ satisfy comparison.  To see this, embed $Z$
as an observable $\widetilde Z$ on the three input ports and equip the
joint observable space with either arithmetic inner product.  The centered
one-port subspaces are mutually orthogonal because the joint state is
$\rho_1\otimes\rho_2\otimes\rho_D$.  The centered slice
$T_iZ-\Tr\gamma Z\,I$ is the orthogonal projection onto the $i$th such
subspace: its pairing with a bounded one-port observable is the defining
partial-expectation identity.  Consequently the sum of the two slice
variances is at most the corresponding variance in $\gamma$.
The variance in $\rho_0=(1-\kappa)\gamma+\kappa\tau_{r_0}$ is at least
$(1-\kappa)$ times that variance in $\gamma$.  Since
$w_i/\lambda_i=1-\kappa$, these statements give
\eqref{eq:weight-comparison} for both multiplication orders.
Only independence across ports has been used.

The two constant families $r_s$ and $r_s+1$ also satisfy comparison.
Indeed, Proposition~\ref{prop:hessian} and \eqref{eq:slack} give
\[
 \sum_i w_i r_i\|L_iZ\|_i^2
 \le (1+\varepsilon)\sum_i\lambda_i r_i\|Z\|_0^2
 \le r_0\|Z\|_0^2,
\]
and the same argument applies with $r_s+1$.

\begin{proposition}[Interpolated comparison]\label{prop:metric}
The positive weights
\begin{equation}\label{eq:thermal-metric}
 F_s(v)=\frac1{h_s}
 \frac{f(h_sv/2)f(-h_s/2)}{f(h_s(v-1)/2)}
 =\frac{b_s(v)-(r_s+\tfrac12)v}{1-v}
\end{equation}
satisfy \eqref{eq:weight-comparison}.  The quotient at $v=1$ is understood
continuously.  For each fixed $r_s>0$ there are constants
$0<c_s\le C_s<\infty$, independent of the minimizing states, such that
$c_s\le F_s(v)\le C_s$ for every $v\in\R$.
\end{proposition}

\begin{proof}
Apply Theorem~\ref{thm:interpolation} with
\[
 m=1/h_s,\qquad x=h_sv_s/2,\qquad y=-h_s/2.
\]
The ambient coordinate spaces in Theorem~\ref{thm:interpolation} are the full weighted matrix spaces; its source test space consists of the centered finite-eigenblock tests described above. Multipliers depending on the frequency preserve that test space, because every diagonal entry has frequency zero. The first two weights are $R_{s,+},R_{s,-}$; the other two are
$r_s,r_s+1$.  The conclusion is the first expression in
\eqref{eq:thermal-metric}.  The second follows from the hyperbolic addition
formula and $\coth(h_s/2)=2r_s+1$.
The first expression is continuous and strictly positive, while the second
gives the limits $r_s$ at $+\infty$ and $r_s+1$ at $-\infty$.
This proves both bounds.
\end{proof}

The thermal weight has a useful logarithmic-mean interpretation:
\begin{equation}\label{eq:tilted-logmean}
 \ell(p_l,p_r)F_s(v_s(l,r))
 =\ell\bigl(r_sp_l,(r_s+1)p_r\bigr).
\end{equation}
Indeed the logarithmic difference on the right is
$h_s(v_s(l,r)-1)$; substitution in \eqref{eq:thermal-metric} proves the
identity, including its continuous values. Up to a fixed normalization,
this is the frequency-dependent logarithmic mean used in the
detailed-balance entropy geometry of Carlen and Maas
\cite[Definition~5.7 and Lemma~5.8]{CarlenMaas2017}.
Proposition~\ref{prop:metric} establishes the particular comparison
needed at our regularized minima.

Write $\|Z\|_{F_s}=\|Z\|_{F_s,s}$, and let
$\|q\|_{F_s,*}$ denote the dual norm of a linear functional on the centered
space.  For a list of $n$ functionals, its squared norm is the sum of the
$n$ squared dual norms.  Density and Proposition~\ref{prop:hessian} extend
Proposition~\ref{prop:metric} to the completed centered spaces.
On bounded observables the extension of $L_i$ is the actual centered
slice.  In fact, bounded spectral compressions of $Z$ converge strongly
with their adjoints, as do their slices.  Dominated convergence in either
arithmetic norm gives convergence in the logarithmic-mean norm, since
$\ell(a,b)\le(a+b)/2$.  This also explains why the comparison applies to
fixed number-basis tests later, although the minimizing eigenbases vary.

\subsection{The thermal defect functionals}
Let $d_{s,j}$ be the annihilator of mode $j$ on port $s$, and put
\[
 \mu_s=(\Tr\rho_s d_{s,j})_{j=1}^n,\qquad
 d'_{s,j}=d_{s,j}-(\mu_s)_jI,\qquad
 e'_s=\Tr\rho_sN_s-|\mu_s|^2.
\]
For bounded $Z$, define
\begin{equation}\label{eq:thermal-defect}
 q_{s,j}(Z)=(r_s+1)\Tr\rho_s(d'_{s,j})^\dagger Z
             -r_s\Tr\rho_sZ(d'_{s,j})^\dagger.
\end{equation}
A coherent displacement shifts $d_{s,j}$ and $(\mu_s)_j$ by the same
constant, so these centered functionals also vanish on displaced thermal
states. The final argument must control $\mu_s$ separately to recover the
fixed, undisplaced thermal reference.
These are well-defined weak trace pairings and vanish on $I$.
For example, $d\rho=(d\rho^{1/2})\rho^{1/2}$ is trace class by
Hilbert--Schmidt factorization whenever $\Tr\rho N<\infty$; the same
holds for $d^\dagger\rho$, $\rho d$, and $\rho d^\dagger$.
Thus no domain invariance for an arbitrary bounded $Z$ is being assumed.

We record the summability needed for both the metric estimate and the
subsequent generator calculation.

\begin{lemma}[Dual norm and logarithmic sums]\label{lem:dual-finiteness}
At each minimizing state, $q_s$ has finite dual norm.  Set
\[
 K_{s,lr}=(r_s+1)p_r-r_sp_l.
\]
Then
\begin{equation}\label{eq:dual-norm-sum}
 \|q_s\|_{F_s,*}^2
 =\sum_{j,l,r}(1-v_s(l,r))K_{s,lr}|(d'_{s,j})_{lr}|^2.
\end{equation}
The sums obtained by replacing the coefficient on the right by
$|K_{s,lr}|$ or $|v_s(l,r)K_{s,lr}|$ are finite.
\end{lemma}

\begin{proof}
Suppress the label $s$ and write $H=-\log\rho$.
By Proposition~\ref{prop:regularized-minimum},
$\Tr\rho W^7<\infty$ and $0\le H\le CW^4$ as forms.
Every eigenvector $\psi_r$ of $\rho$ belongs to $D(W^{7/2})$, since
$p_r>0$.  Number shifts give
$\|W^2d'\psi\|+\|W^2(d')^\dagger\psi\|
 \le C'\|W^{5/2}\psi\|$.
Consequently the crossed logarithmic sums, such as
\[
 \sum_{l,r}p_r(-\log p_l)|d'_{lr}|^2
 =\sum_r p_r\langle d'\psi_r,Hd'\psi_r\rangle,
\]
are bounded by a constant times $\Tr\rho W^5$.
For the uncrossed sums, use
\[
 -\log p_r\le C\langle\psi_r,W^4\psi_r\rangle,
 \qquad
 \langle W^4\rangle_{\psi_r}\langle W\rangle_{\psi_r}
 \le\langle W^5\rangle_{\psi_r}.
\]
The last inequality is the elementary nonnegative covariance inequality
for the increasing functions $t^4,t$ of the same positive random variable.
It bounds $(-\log p_r)\|d'\psi_r\|^2$ and its creation-operator analogue.
The corresponding sums without logarithms are finite by finite energy.
These bounds imply the claimed absolute summability of $K$ and $vK$.

In matrix coordinates, \eqref{eq:thermal-defect} is
\[
 q_j(Z)=\sum_{l,r}K_{lr}\overline{d'_{j,lr}}Z_{lr}.
\]
By \eqref{eq:arithmetic-weights} and \eqref{eq:thermal-metric},
\begin{equation}\label{eq:defect-multiplier}
 \frac{K_{lr}}{h\ell(p_l,p_r)}
 =b(v)-(r+\tfrac12)v=(1-v)F(v).
\end{equation}
Hence the representing vector in the $F$ inner product is
$(1-v)d'_j$, and its squared norm is the sum in
\eqref{eq:dual-norm-sum}.  That sum is finite by the preceding estimates.
Finally, $F(0)=1/h$, so orthogonality to $I$ in the $F$ inner product is
exactly zero $\rho$-mean.  The representing vector is orthogonal to $I$
because $q_j(I)=0$.  Restricting to the centered space therefore does not
change the dual norm.
\end{proof}

\begin{lemma}[Weak convolution identity]\label{lem:weak-traces}
Let $\mu_\gamma=\sum_i\sqrt{\lambda_i}\mu_i$ be the raw output mean.
For every bounded output observable $Z$,
\begin{equation}\label{eq:defect-convolution}
 q_{0,j}(Z)=\sum_i w_iq_{i,j}(L_iZ)
 +\kappa(1-\kappa)\overline{(\mu_\gamma)_j}
                  \Tr[(\gamma-\tau_{r_0})Z].
\end{equation}
\end{lemma}

\begin{proof}
The thermal defect is zero on $\tau_r$, since
$(r+1)d\tau_r=r\tau_rd$ by its geometric number weights.
For the raw output, define $q_{\gamma,j}$ by
\eqref{eq:thermal-defect} using its mean and reference parameter $r_0$.
The terms with coefficient $1$ in the identity
\[
 q_{\gamma,j}(Z)
 =\sum_{s=1,2,D}\sqrt{\lambda_s}\,q_{s,j}(T_sZ)
\]
agree by $d_{0,j}=\sum_s\sqrt{\lambda_s}d_{s,j}$.
The remaining terms are expected commutators.  After rotation,
$d_{s,j}^\dagger$ is $\sqrt{\lambda_s}d_{0,j}^\dagger$ plus an operator
on discarded modes.  The weak commutator of that discarded operator with
a retained observable has zero trace.  Thus its retained coefficient is
$\sum_s\lambda_sr_s=r_0$, as required.

These statements use only finite energy: all one-mode-factor products
with the state are trace class as above.  A discarded operator can first
be cut off in its own number basis; its bounded cutoff commutes with
$Z\otimes I$, so trace cyclicity applies.  The cutoff products converge
in trace norm by the same Hilbert--Schmidt factorization.  This justifies
the weak commutator and passage between joint, reduced, and sliced traces,
without assigning an operator domain to $[d_{0,j}^\dagger,Z]$.

The $D$ term vanishes.  Since the mean of the replaced output is
$(1-\kappa)\mu_\gamma$, direct centering gives
\[
 q_{0,j}=(1-\kappa)q_{\gamma,j}
  +\kappa(1-\kappa)\overline{(\mu_\gamma)_j}
       \Tr[(\gamma-\tau_{r_0})\,\cdot\,].
\]
The first term is $\sum_iw_iq_{i,j}\circ L_i$, since $q_{i,j}(I)=0$.
\end{proof}

\begin{proposition}[Defect estimate]\label{prop:defect-estimate}
Set
\[
 A_*^2=\sum_iw_i\|q_i\|_{F_i,*}^2,
 \qquad M_*^2=\sum_iw_i|\mu_i|^2,
 \qquad C_0=(h_0\inf_{v\in\R}F_0(v))^{-1/2}.
\]
Then
\begin{equation}\label{eq:defect-triangle}
 \|q_0\|_{F_0,*}\le A_*+\sqrt{2\kappa}\,C_0M_*.
\end{equation}
In particular, the previously reserved choice of $\kappa$ may be made
small enough that
\begin{equation}\label{eq:defect-estimate}
 \|q_0\|_{F_0,*}^2
 \le A_*^2+\frac\varepsilon2(A_*^2+M_*^2)
\end{equation}
holds uniformly in $\zeta$ and the minimizing states.
\end{proposition}

\begin{proof}
By Proposition~\ref{prop:metric} and Cauchy--Schwarz, the dual norm of the
first term in \eqref{eq:defect-convolution}, summed over the modes, is at
most $A_*$.  For the second term put
\[
 J(Z)=\sqrt{\kappa(1-\kappa)}\Tr[(\gamma-\tau_{r_0})Z].
\]
The variance between the two branches of the mixture gives
$|J(Z)|^2\le\|Z\|_{R_{0,+},0}^2$ and
$|J(Z)|^2\le\|Z\|_{R_{0,-},0}^2$ on centered tests.
The two-weight logarithmic-mean operation proved in
Section~\ref{int:operations} applies to this scalar-valued map,
with both target weights equal to $1$.
Since $\ell(R_{0,+},R_{0,-})=1/h_0$, it yields
\[
 |J(Z)|^2\le\|Z\|_0^2/h_0
 \le C_0^2\|Z\|_{F_0}^2.
\]
Moreover,
$(1-\kappa)|\mu_\gamma|^2\le2\sum_iw_i|\mu_i|^2=2M_*^2$.
The triangle inequality now proves \eqref{eq:defect-triangle}.
Using Young's inequality with coefficient $\varepsilon/2$, it is enough
to impose
\[
 0<\kappa\le\frac{\varepsilon^2}{4(\varepsilon+2)C_0^2}
\]
to obtain \eqref{eq:defect-estimate}.  The right-hand side depends only on
the reference parameters and $\varepsilon$, which were fixed before
$\kappa$; this requirement is therefore compatible with the parameter
choice in Proposition~\ref{prop:regularized-minimum}.
\end{proof}
